% !TeX root = ../practica.tex
% !TeX encoding = utf8
%
% ====================================================================
%  BIBLIOGRAFÍA Y FUENTES CONSULTADAS T6, 7 y 8
% ====================================================================

\chapter*{Fuentes consultadas (Temas 6, 7 y 8)}
\addcontentsline{toc}{chapter}{Fuentes consultadas (Temas 6, 7 y 8)}

% ================================================================
\section*{Fuentes sobre CaixaBankCareers y procesos de Selección}
% ================================================================

\begin{enumerate}[label={[\arabic*]}]

  \item \textbf{CaixaBankCareers} --- ``Página oficial de la bolsa de empleo e información de etapas de reclutamiento''.\\
        \url{https://caixabankcareers.com/}

  \item \textbf{Jobseeker} --- Análisis de los requisitos para la criba curricular.\\
        \url{https://www.jobseeker.com/es/}

  \item \textbf{Ejemplos-Curriculum} --- Revisión sobre los procesos bancarios automatizados y filtros (\textit{Killer questions}).\\
        \url{https://www.ejemplos-curriculum.com/blog/entrevista-trabajo-la-caixa/}

  \item \textbf{Foros de empleo y Oposiciones} (Infoempleo, Randstad, Buscaoposiciones) --- Experiencias reales de candidatos en dinámicas de grupo y pruebas psicotécnicas de CaixaBank.\\
        \url{https://www.infoempleo.com/}
% Fuentes de jose angel:
  \item \textbf{LinkedIn --- CaixaBank} --- Perfil corporativo; análisis de la estrategia de \textit{employer branding} y publicación de ofertas (reclutamiento 3.0).\\
        \url{https://es.linkedin.com/company/caixabank/jobs}

  \item \textbf{People Xperience Hub / Dialnet} --- Artículo académico sobre la iniciativa RPO de CaixaBank (2019): estandarización y digitalización del ciclo de atracción de talento.\\
        \url{https://dialnet.unirioja.es/servlet/articulo?codigo=7731413}

  \item \textbf{Entrevista con directora de sucursal de CaixaBank} --- Entrevista semiestructurada realizada el 11 de abril de 2026. Fuente primaria sobre el proceso de reclutamiento y selección aplicado en la entidad: métodos utilizados, herramientas digitales, videoentrevistas e iniciativas en curso (IA, gamificación).

\end{enumerate}

% ================================================================
\section*{Documentos teóricos de referencia}
% ================================================================

\begin{enumerate}[label={[\arabic*]}, resume]

  \item Apuntes de clase de los Temas 6, 7 y 8.

  \item \textbf{Guión para el desarrollo de la práctica grupal Temas 6, 7 y 8},
        Dirección de Recursos Humanos~I, Doble Grado ADE-Ingenierías,
        Universidad de Granada, curso 2025--2026.

\end{enumerate}

\endinput
